\chapter{검증 결과(7): Stage A 경부하 운전과 모듈 관리·고장}

\section{위상 셰딩의 필요성과 설계}

영교차 추종 게이트는 전류가 검출 문턱(0.5 A) 아래로 돌아온 뒤 검출 지연(20 ns)만큼 늦게 끄므로\cite{b22}, 공진전류 첨두가 작으면 턴오프 전류/피크가 오른다. 같은 출력을 위상 하나가 맡으면 첨두가 두 배가 되어 비율이 절반 아래로 내려가고 위상당 고정 손실도 준다\cite{b23}. 운전점 계산(E24, 1008 V)에서 2위상의 턴오프 전류/피크는 모듈 출력 10 A에서 5.51 \%로 기준 5 \%를 넘었고, 위상 하나는 1.58 \%였다. 충전 정격점(788 V, 위상 전류 46.9 A)의 위상 하나 첨두 76.4 A를 위상 정격 첨두로 둔다. 11.1–11.4절의 시험은 150 kHz Stage A로, 모듈 손실은 출력용량 손실을 정정하기 전 모델로 했고(13장), 최종 설계(A25R)의 문턱과 판정은 11.5절에 있다.


\begin{itemize}
\item \textbf{멈춤 15 A(100 µs 지속), 복귀 22 A 또는 남은 위상의 순시 전류 36 A.} 2위상은 12 A 이상에서 기준 5 \% 안이므로 3 A 여유를 두었고, 필터한 평균전류는 부하 계단을 늦게 보므로 순시 전류도 본다. 멈출 때는 도통 중인 구간을 자르지 않아 멈추는 순간도 영전류이며, 복귀는 바로 한다.
\end{itemize}

확인 과정에서 Stage A 스위칭 모델의 데드타임 처리 오차를 찾아 고쳤다(다이오드 환류 전류가 적분 간격 안에 0으로 가지 않아 진동). 회로(TB16·TB17) V\textsubscript{A}와의 최대 차가 0.094 V에서 0.014 V로 줄었다. 이전 처리는 턴오프 비율을 보수 쪽으로 계산했으므로 E1–E23의 판정은 바뀌지 않고, E24 이후의 기준값은 고친 처리로 계산했다.

\section{위상 셰딩의 결과}

셰딩은 새 과도(멈추는 순간, 위상 하나의 부하 계단, 재시작, 문턱 근처 반복 전환)를 만든다. 그래서 정상 운전점과 전환 과도를 시험 전에 정한 기준(모든 턴오프 |i\textsubscript{off}| ≤ max(5 \% × 반주기 첨두, 1.0 A), 위상 첨두 ≤ 76.4 A, 반복 전환 없음)으로 판정했다(표 \ref{t:e24}, 그림 \ref{f:e24}).

{\footnotesize
\begin{longtable}{L{\dimexpr 0.2885\linewidth-2\tabcolsep\relax}L{\dimexpr 0.5962\linewidth-2\tabcolsep\relax}L{\dimexpr 0.1154\linewidth-2\tabcolsep\relax}}
\caption{Stage A 경부하 위상 셰딩의 판정(E24, TB18, TB19)}\label{t:e24}\\
\toprule
\tabhead \textbf{시험} & \textbf{결과} & \textbf{판정} \\
\midrule
\endfirsthead
\toprule
\tabhead \textbf{시험} & \textbf{결과} & \textbf{판정} \\
\midrule
\endhead
\bottomrule
\endlastfoot
E24 정상 운전점(1008·540 V × −10–75 A) & 7.5 A 이상 턴오프 전류/피크 ≤ 3.27 \%(셰딩 없는 10 A 5.51 \%). 7.5 A 미만은 턴오프 전류 ≤ 0.48 A & 26/26 \\
E24 과도(계단·램프·역방향·충전기 주입 10사례) & 턴오프 위반 0, 반복 전환 없음. V\textsubscript{A} 노드 직접 계단에서 위상 첨두 72.9 A(셰딩 없는 60.5 A). 사례 E(모듈당 충전 램프 요구의 4배)는 셰딩 없이도 82.5 A, 셰딩 93.6 A로 정격 초과 & 9/10 \\
TB18 회로(R2 버스 뒤 부하 계단·램프, 충전기 주입, 역방향 5사례) & 턴오프 3,547번 위반 0(최대 0.77 A), 위상 첨두 ≤ 52.5 A. E24와 첨두 0.23 A 안 & 5/5 \\
TB19 회로(TB17 + 셰딩 게이트) & 1008 V 강등 뒤 10 A의 턴오프 전류/피크 4.44 → 1.31 \%. 턴오프 17,561번 위반 0, 위상 첨두 ≤ 48.6 A. 강등·48 V·Stage C 결과는 TB17과 같음 & 2/2 \\
\end{longtable}
}
\srcnote{주: 7.5 A 미만에서는 턴오프 전류/피크 비율이 뜻이 없어 절대값(≤ 1.0 A, 검출 문턱의 두 배)으로 판정했다.}

\begin{figure}[!htb]
\centering
\includegraphics[width=15.5cm]{\repo/validation/2026-10-26-stage-a-shedding/E24/E24_stage_a_shedding.png}
\caption{Stage A 경부하 위상 셰딩(E24)}\label{f:e24}
\srcnote{주: 150 kHz 설계의 결과이며 (b)는 출력용량 손실을 정정하기 전의 모델 값이다(최종 설계의 문턱과 판정은 표 \ref{t:a25r_mgr}). (a) 모듈 출력전류별 턴오프 전류/피크, (b) 모듈 손실, (c) 54 → 10 → 54 A 계단의 위상 전류와 셰딩 모드, (d) 사례 E(모듈당 충전 램프 요구의 4배)의 위상 전류.}
\end{figure}

불통과 1건(사례 E)은 셰딩이 없어도 위상 정격을 넘는 조건이다. 셰딩은 복귀 지연만큼 첨두를 11.1 A 더했다. 150 kHz 설계의 램프 요구(전체 3.75 A/µs, 모듈당 0.94 A/µs, 8.3절)에서는 위상 첨두가 52.6 A로 셰딩 없는 경우(53.1 A)와 같아, 충전기 쪽 램프 제한이 위상 정격도 지킨다. 

\section{4모듈 관리기와 부품 공차(E28, TB22)}

운전점마다 필요한 모듈 수가 다르다(주행 수–100 A: 1–2개, 충전 375 A: 4개, R2-W460 강등 뒤 약 10 A: 1개의 위상 하나). 모듈 넷을 늘 돌리면 경부하 영전류 스위칭 비율이 나빠지고 고정 손실이 여덟 위상분 들므로, 동작 모듈 수를 정하는 관리기를 설계했다.

\begin{itemize}
\item \textbf{추가(바로).} 전체 출력전류 > 동작 수 × 70 A, 위상 순시 전류 > 66 A, 충전 요청 가운데 하나면 모듈을 더한다.
\item \textbf{줄임(지속 뒤).} ⌈전체 전류/60 A⌉ < 동작 수가 1 ms 이어지면 줄이고, 그 모듈은 지금 구간을 자연 턴오프로 끝내고 멈춘다. 문턱 60·120·180 A는 모듈 N개와 N+1개의 손실 교차점보다 낮다.
\item \textbf{충전 시작.} 충전 요청이 오면 충전기 전류가 오르기 전에 모듈을 모두 켠다(사전 기동).
\item \textbf{동작 모듈.} 모듈마다 11.1절의 위상 셰딩을 한다.
\end{itemize}

스위칭 모델(4모듈 8위상, R2 버스, V\textsubscript{A} 노드 충전기)로 공칭과 공차 표본 3개(C 5 \%, L 10 \%, R 10 \%)를 계산했다. 1·2차 실행의 불통과로 설계를 네 군데 고친 뒤 시험 전에 정한 기준으로 3차 실행을 판정했다. \textbf{정상 운전점은 공칭 33/36, 공차 99/108}(턴오프 비율 ≤ 4.1 \%, 턴오프 전류 ≤ 0.91 A)이었고, 불통과는 정격 충전점(375 A)의 첨두뿐으로 판정 창에 공진전류 0에서 시작하는 시작 과도(82.1 A)가 들어가기 때문이다(판정 창은 바꾸지 않음). \textbf{과도는 공칭 6/7, 공차 21/21}로, 불통과는 사전 기동이 없는 충전(위상 첨두 84.2 A)이어서 사전 기동을 요구사항에 넣었다. 회로(TB22)는 사전 기동이 있는 세 사례 통과, 없는 사례 불통과로 모델과 같았다(전환 시각 2 µs, 위상 첨두 3.3 A 안). 공차 표본의 정격 충전점 위상 첨두는 공칭보다 10–12 \% 컸고 위상 평균전류 불균형은 최대 8.5 \%로, 부품 정격 표(16장)의 입력이 되었다. \textbf{멈춘 모듈의 플라잉 커패시터 전압}은 V\textsubscript{bat}/2 대비 −10–0 V일 때만 재시작이 판정을 지켰다(+2 V에서 남은 모듈 첨두 77.1 A, −15 V에서 재시작 첨두 106.9 A, 150 kHz 설계에서만 확인). \textbf{게이트 변경.} 무장하지 못한 반주기(첨두 < 2.5 A)는 부호 없이 |i| < 0.5 A가 검출 지연만큼 이어질 때 끄도록 바꿨고(12.3절의 방향 반전 처리를 대체), E27 사례 55개와 E24 판정을 다시 확인해 결과가 같았다.

\section{모듈 고장 처리(E30, TB24)}

E28의 관리기는 고장 모듈을 후보에서 빼지 않아, 주행 중 유일한 동작 모듈이 고장 나면 대신할 모듈이 시작하지 않았고 충전 중에는 남은 3모듈이 수십 ms 이상 375 A를 받았다. 그래서 고장 간격에 고장 모듈 위상의 네 스위치를 끄고(남은 공진전류는 바디 다이오드로 소멸), 건전 모듈 목록에서 영구히 빼며, 동작 모듈이 0이 되면 다음 건전 모듈을 바로 두 위상으로 시작하게 했다.

주행·충전 중 고장 8사례를 공칭·공차 표본으로 판정했다(기준: 턴오프 |i\textsubscript{off}| ≤ max(5 \% × 반주기 첨두, 1.0 A), 위상 첨두 주행 ≤ P\textsubscript{rated}·충전 ≤ 1.5 × P\textsubscript{rated}, V\textsubscript{A} 고장 직전 ±10 V, 끝 동작 수). \textbf{모델 32/32, 회로(TB24) 4/4로 통과했다.} 주행 중 대체 지연은 0 µs, V\textsubscript{A} 변화는 0.2 V 안이었다. 충전 중 남은 3모듈의 위상 첨두는 103.5–114.5 A(한도 114.7–128.3 A)로, 이 값이 Stage A 소자·인덕터·플라잉 커패시터의 정격 조건이 되었다(16장). 회로는 모델과 위상 첨두 0.9 A, V\textsubscript{A} 0.05 V 안이었다.

\section{최종 설계(A25R)의 문턱과 판정}

A25R로 바꾸면서(13장) 문턱과 시간 상수를 다시 정하고 같은 기준으로 다시 확인했다(E32, E35, TB27, 표 \ref{t:a25r_mgr}).

{\footnotesize
\begin{longtable}{L{\dimexpr 0.26\linewidth-2\tabcolsep\relax}L{\dimexpr 0.30\linewidth-2\tabcolsep\relax}L{\dimexpr 0.44\linewidth-2\tabcolsep\relax}}
\caption{Stage A 경부하·모듈 관리의 150 kHz 설계 값과 최종 설계 값}\label{t:a25r_mgr}\\
\toprule
\tabhead \textbf{항목} & \textbf{150 kHz 설계} & \textbf{A25R(최종)} \\
\midrule
\endfirsthead
\toprule
\tabhead \textbf{항목} & \textbf{150 kHz 설계} & \textbf{A25R(최종)} \\
\midrule
\endhead
\bottomrule
\endlastfoot
위상 셰딩(멈춤 / 복귀 / 순시) & 15 / 22 / 36 A & 20 / 30 / 53 A(두 위상 7.5–12.5 A의 턴오프 비율 때문) \\
모듈 관리기(추가 / 줄임 / 순시) & 70 / 60 A/모듈, 66 A & 같음(손실 교차 89·156·221 A ≥ 줄임 60·120·180 A) \\
시간 상수(필터, 멈춤 지속, 모듈 필터, 잠금, 유지) & 2 µs, 100 µs, 20 µs, 20 µs, 100 µs & 6배: 12 µs, 600 µs, 120 µs, 120 µs, 600 µs(모듈 줄임 지속 1 ms는 그대로) \\
정격 충전점 위상 첨두 P\textsubscript{rated} & 76.4–76.5 A & 74.5 A(공차 표본 86.5–102.8 A) \\
충전 전 사전 기동·충전전류 램프 & ≥ 100 µs & ≥ 600 µs \\
정상 운전점(공칭 / 공차 표본) & 33/36 / 99/108 & 33/36 / 36/36(공칭 불통과는 정격 충전점 시작 과도 76.7 A, 마지막 10주기 74.3 A) \\
요구 안 과도(공칭 / 표본) & 6/7 / 21/21 & 6/6 / 6/6 \\
모듈 고장(검출 즉시 해당 모듈 게이트 차단) & 32/32 & 16/16(3모듈 충전 첨두 최대 105.5 A; 후보 C\textsubscript{fly} 45.0·42.75 µF에서 106.86 A) \\
회로 & TB22 3/4, TB24 4/4 & TB27 4/4(첨두 차 ≤ 2.2 A, V\textsubscript{A} 차 ≤ 0.06 V), TB28 4/4 \\
\end{longtable}
}
\srcnote{주: A25R의 셰딩 문턱 20 / 30 / 53 A는 지정한 모듈·부하 이력에서 정한 제어값이다. 충전 시작은 보조부하 15·25 kW에서 확인했고, 640 V·보조부하 1 kW에서는 같은 기동·램프로 78.74 A가 되어 공칭 한도 74.53 A를 넘었다(E46). 150 kHz 설계에서 확인한 프리차지·정지 모듈 재기동은 최종 C\textsubscript{A} 2.4 mF의 냉기동 근거로 쓰지 않는다(최종 설계로 다시 확인하지 않음). 모듈 고장은 고장 표지를 받는 즉시 해당 모듈의 모든 게이트를 끄는 차단 가능한 고장이며, 강제 차단 순간은 정상 턴오프 평가와 분리했다(단락 소자·게이트 고착·검출 지연은 미확인). 부하 계단은 요구사항의 계약(Tier-1 대표 시험 ±10 kW)으로 판정했다. 150 kHz부터 써 온 20 kW 순간 계단은 계약 밖으로 정보로만 계산했고 A25R에서 불통과였다(첨두 87–88 A). 자료: 검증 근거 E24, E28, E30, E32, E35.}

시간 상수를 6배로 늘린 것은 공진 주기가 6배 길어졌기 때문이다. 셰딩 문턱을 올린 것은 A25R의 두 위상 운전이 모듈 출력 7.5–12.5 A에서 턴오프 비율 8.7–5.0 \%로 기준(5 \%)을 넘었기 때문이다(위상 하나는 7.5 A 이상에서 4.1 \% 이하, E32). 멈춘 모듈의 플라잉 커패시터 전압 요구는 A25R에서 다시 확인하지 않았다(18장).
